% ======================================================================
%  MÓDULO 11 — Referências auditadas e metodologia de citação
%  Regra R110: DOI citado deve ser ativo e verificado externamente.
% ======================================================================
\thispagestyle{fancy}
\chapter{Referências Auditadas e Metodologia de Citação}

\roteirodidatico
{Uma citação permite ao leitor localizar a fonte usada em uma afirmação. Para servir a esse propósito, a referência precisa estar completa e a fonte precisa sustentar o trecho associado.}
{\emph{DOI} identifica um objeto digital; \emph{citação direta} reproduz palavras da fonte; \emph{paráfrase} reformula a ideia; \emph{ABNT} define padrões bibliográficos adotados no manuscrito.}
{Confira autor, título, veículo, ano e localização; abra a fonte; confronte texto original, tradução e afirmação; registre data e método de acesso.}
{Um DOI resolvível confirma a localização do registro, não a interpretação nem a pertinência da citação. Traduções são identificadas como livres e alegações devem permanecer dentro do escopo da fonte.}
{Um leitor consegue encontrar a fonte e verificar, sem depender da interpretação do autor, a afirmação citada e sua localização?}

\casoguiado{conferir uma citação}
{Diante de ``(SOBRENOME, ano, p. 12)'', localize a referência completa, abra a edição citada e confira a página. Se for uma paráfrase, compare a ideia e preserve suas qualificações; se for citação direta, confira cada palavra e o localizador. Um DOI pode levar ao artigo, mas não necessariamente à página citada ou ao texto integral.}

\begin{resultado}
Toda citação com DOI deste manual segue a regra \textbf{R110}/\textbf{anti-overclaim}
(\emph{AGENTS.md}, item 4): nenhum DOI é citado ``de memória''. Antes de entrar
no texto, a referência é \textbf{verificada por consulta externa} (webfetch do
resumo/página do periódico ou do repositório da editora), o \textbf{trecho
original} é copiado literalmente do documento e sua \textbf{tradução livre}
para o português é rotulada como tal. Um DOI que não resolve em
\pth{https://doi.org} ou cujo trecho não confere com a fonte \textbf{não entra}
no livro. Verificação corrente: \textbf{29 de setembro de 2026}.

\textbf{Normas ABNT atualizadas.} As referências seguem a
\textbf{ABNT NBR 6023:2018} (Informação e documentação --- Referências:
elaboração) e as citações seguem a \textbf{ABNT NBR 10520:2023} (Informação e
documentação --- Citações em documentos: apresentação). Adota-se o
\textbf{sistema numérico}: a chamada no texto é a remissão à nota de rodapé, que
traz a \textbf{referência completa} em ABNT, o DOI com \textbf{data de acesso}
e, quando há citação direta, o \textbf{trecho textual} e sua página.
\end{resultado}

\section{Biblioteca de referências auditadas}

\begin{referencia}
\textbf{FLAVELL, J. H.} Metacognition and cognitive monitoring: A new area of
cognitive-developmental inquiry. \emph{American Psychologist}, v.~34, n.~10,
p.~906-911, 1979. DOI: \href{https://doi.org/10.1037/0003-066X.34.10.906}{10.1037/0003-066X.34.10.906}. Acesso em: 29 set. 2026.
\textbf{Uso no manual:} Módulo 3 (definição de conhecimento metacognitivo
persistido no MetaBus). \textbf{Status:} DOI ativo; resumo conferido na fonte
original (Stanford/Sage, out.~1979).
\end{referencia}

\begin{referencia}
\textbf{NII, H. P.} Blackboard systems: The blackboard model of problem solving
and the evolution of blackboard architectures. \emph{AI Magazine}, v.~7, n.~2,
1986. DOI: \href{https://doi.org/10.1609/aimag.v7i2.537}{10.1609/aimag.v7i2.537}. Acesso em: 29 set. 2026.
\textbf{Uso no manual:} Módulo 3 (proveniência do padrão blackboard/A2A e do
vocabulário quadro--fonte de conhecimento). \textbf{Status:} DOI ativo; resumo
conferido no AAIAI Magazine; relato do HEARSAY-II (1971--1976) confirmado.
\end{referencia}

\begin{referencia}
\textbf{IOANNIDIS, J. P. A.} Why most published research findings are false.
\emph{PLoS Medicine}, v.~2, n.~8, e124, 2005. DOI:
\href{https://doi.org/10.1371/journal.pmed.0020124}{10.1371/journal.pmed.0020124}. Acesso em: 29 set. 2026.
\textbf{Uso no manual:} Módulos 3 e 6 (fundamento empírico do gate
\cod{no_verified_without_external} e da regra anti-overclaim).
\textbf{Status:} DOI ativo; versão integral acessível em acesso aberto
(PubMed Central PMC1182327); simulações e PPV conferidos no texto integral.
\end{referencia}

\begin{referencia}
\textbf{WILKINSON, M. D.; DUMONTIER, M.; AALBERSBERG, I. J. et al.} The FAIR
Guiding Principles for scientific data management and stewardship.
\emph{Scientific Data}, v.~3, 160018, 2016. DOI:
\href{https://doi.org/10.1038/sdata.2016.18}{10.1038/sdata.2016.18}. Acesso em: 29 set. 2026.
\textbf{Uso no manual:} Módulos 4 e 5 (proveniência defensável dos registros de
execução e reprodutibilidade do patrimônio de specs/testes).
\textbf{Status:} DOI ativo; princípios F/A/I/R (incl. R1.2 proveniência) e
licença CC BY 4.0 conferidos no artigo original (Nature Portfolio).
\end{referencia}

\begin{referencia}
\textbf{WOOLDRIDGE, M.; JENNINGS, N. R.} Intelligent agents: theory and
practice. \emph{The Knowledge Engineering Review}, v.~10, n.~2, p.~115-152,
1995. DOI:
\href{https://doi.org/10.1017/S0269888900008122}{10.1017/S0269888900008122}. Acesso em: 29 set. 2026.
(4093 citações indexadas no Crossref em 29/09/2026; resumo conferido.)
\end{referencia}

\begin{referencia}
\textbf{JENNINGS, N. R.; SYCARA, K.; WOOLDRIDGE, M.} A roadmap of agent
research and development. \emph{Autonomous Agents and Multi-Agent Systems},
v.~1, n.~1, p.~7-38, 1998. DOI:
\href{https://doi.org/10.1023/A:1010090405266}{10.1023/A:1010090405266}. Acesso em: 29 set. 2026.
(Objetivo conforme resumo em Semantic Scholar: identificar conceitos-chave e
aplicações e sua relação mútua.)
\end{referencia}

\begin{referencia}
\textbf{PENG, R. D.} Reproducible research in computational science.
\emph{Science}, v.~334, n.~6060, p.~1226-1227, 2011. DOI:
\href{https://doi.org/10.1126/science.1213847}{10.1126/science.1213847}. Acesso em: 29 set. 2026.
(1157 citações indexadas no Crossref em 29/09/2026; resumo conferido.)
\end{referencia}

\begin{referencia}
\textbf{SMITH, R. G.} The contract net protocol: high-level communication and
control in a distributed problem solver. \emph{IEEE Transactions on Computers},
v.~C-29, n.~12, p.~1104-1113, 1980. DOI:
\href{https://doi.org/10.1109/TC.1980.1675516}{10.1109/TC.1980.1675516}. Acesso em: 29 set. 2026.
\textbf{Uso no manual:} Módulos 2, 7 e 14 (gerência/contratante do Blackboard e
das bridges de integração). \textbf{Status:} DOI ativo; resumo conferido em
\pth{reidgsmith.com}; a licitação e a negociação de contrato conferidas no
texto integral disponível. (Não confundir com o DOI
10.1109/TC.1980.1675514, que é o artigo de Gonzalez \& Jordan sobre avaliação
quantitativa de sistemas distribuídos.)
\end{referencia}

\begin{referencia}
\textbf{BROOKS, F. P. Jr.} No silver bullet: essence and accidents of software
engineering. \emph{Computer}, v.~20, n.~4, p.~10-19, 1987. DOI:
\href{https://doi.org/10.1109/MC.1987.1663532}{10.1109/MC.1987.1663532}. Acesso em: 29 set. 2026.
\textbf{Uso no manual:} Módulo 4 (justificativa do protocolo SDD/TDD como
disciplina acumulada, não bala de prata). \textbf{Status:} DOI ativo; 1804
citações indexadas no Crossref em 29/09/2026; trecho de p.~11 conferido no
texto integral (cópias abertas do \emph{Computer}).
\end{referencia}

\begin{referencia}
\textbf{PAGE, M. J. et al.} PRISMA 2020 explanation and elaboration: updated
guidance and exemplars for reporting systematic reviews. \emph{BMJ}, v.~372,
n160, 2021. DOI:
\href{https://doi.org/10.1136/bmj.n160}{10.1136/bmj.n160}. Acesso em: 29 set. 2026.
\textbf{Uso no manual:} Módulos 6, 11 e 14 (relato transparente de revisões e
da biblioteca de citações). \textbf{Status:} DOI ativo; 9559 citações
indexadas no Crossref em 29/09/2026; resumo conferido na página do BMJ.
(Companion do relato de princípios: DOI 10.1136/bmj.n71.)
\end{referencia}

\begin{referencia}
\textbf{AMDAHL, G. M.} Validity of the single processor approach to achieving
large scale computing capabilities. In: \emph{AFIPS Conference Proceedings},
v.~30, p.~483-485, 1967. DOI:
\href{https://doi.org/10.1145/1465482.1465560}{10.1145/1465482.1465560}. Acesso em: 29 set. 2026.
\textbf{Uso no manual:} Módulos 5 e 14 (limite da aceleração paralela do
pipeline/roteamento). \textbf{Status:} DOI ativo; trecho de p.~484 conferido
no fac-símile do texto integral (arquivo CMU); citação canônica da Lei de
Amdahl.
\end{referencia}

\begin{referencia}
\textbf{SIMON, H. A.} A behavioral model of rational choice. \emph{The
Quarterly Journal of Economics}, v.~69, n.~1, p.~99-118, 1955. DOI:
\href{https://doi.org/10.2307/1884852}{10.2307/1884852}. Acesso em: 29 set. 2026.
\textbf{Uso no manual:} Módulos 8 e 14 (racionalidade limitada como base da
decisão assistida por agentes). \textbf{Status:} DOI ativo; 10407 citações
indexadas no Crossref em 29/09/2026; trecho de p.~99 conferido no texto
integral (JSTOR, acesso legitimado).
\end{referencia}

\begin{referencia}
\textbf{NEWELL, A.} The knowledge level. \emph{Artificial Intelligence},
v.~18, n.~1, p.~87-127, 1982. DOI:
\href{https://doi.org/10.1016/0004-3702(82)90012-1}{10.1016/0004-3702(82)90012-1}. Acesso em: 29 set. 2026.
\textbf{Uso no manual:} Módulo 14 (o agente descrito no nível do conhecimento,
acima do nível simbólico). \textbf{Status:} DOI ativo; trecho do resumo
(p.~87) conferido na página do periódico. (Nota: o sufixo correto é
\cod{90012-1}; \cod{90012-2} não resolve.)
\end{referencia}

\begin{referencia}
\textbf{NELSON, T. O.; NARENS, L.} Metamemory: a theoretical framework and new
findings. In: BOWER, G. H. (ed.). \emph{The Psychology of Learning and
Motivation}. Academic Press, v.~26, p.~125-173, 1990. DOI:
\href{https://doi.org/10.1016/S0079-7421(08)60053-5}{10.1016/S0079-7421(08)60053-5}. Acesso em: 29 set. 2026.
\textbf{Uso no manual:} Módulo 3 (monitoramento e controle do MetaBus).
\textbf{Status:} DOI ativo após auditoria da sessão R613; recorte do resumo do
capítulo (p.~125) conferido na página da editora; é capítulo de livro — não
publicar como artigo de periódico.
\end{referencia}

\section{Metodologia de verificação}

\begin{metodo}
\emph{Como uma citação entra neste livro (regra R110):}
\begin{enumerate}[leftmargin=1.4em,itemsep=1pt]
  \item \textbf{Sondagem do DOI} — o link \pth{https://doi.org/<id>} é resolvido
        e a página do periódico/editora é lida (não apenas o resumo agregador).
  \item \textbf{Recorte: resumo ou texto integral} — quando o trecho citado
        vem apenas do resumo (páginas de editora que impedem acesso ao texto
        integral, ex.: BMJ e ACM), o recorte é rotulado ``resumo''
        \emph{e} a página indicada é a página inicial do artigo; quando o
        texto integral está aberto (ex.: Amdahl, Simon, Brooks), o recorte é
        conferido no próprio documento e a página exata do trecho é citada.
        Nunca se cita página de resumo como se fosse página de texto integral.
  \item \textbf{Conferência do trecho} — o recorte citado é copiado
        \emph{literalmente} do documento; citações parafraseadas são marcadas
        como paráfrase, não como trecho.
  \item \textbf{Relevância} — a nota de rodapé declara \emph{por que} aquela
        fonte é pertinente ao contexto daquela ficha (não citação decorativa).
  \item \textbf{Tradução} — a tradução pt-BR é rotulada como tradução livre e
        fiel, mantendo números, negações e entidades do original.
  \item \textbf{Forma (ABNT)} — a referência é redigida conforme a
        \textbf{NBR 6023:2018} (sobrenome em maiúsculas, título da publicação
        em destaque, volume, número, páginas com hífen, ano, DOI e data de
        acesso) e a citação conforme a \textbf{NBR 10520:2023} no
        \textbf{sistema numérico}: a chamada remete à nota de rodapé, que traz
        a referência completa e, em citação direta, o trecho e a página.
  \item \textbf{Atesto} — o DOI permanece na biblioteca apenas se ativo no
        momento da verificação; DOIs que quebram entre edições são removidos
        na próxima auditoria do manual.
\end{enumerate}
\end{metodo}

\begin{aviso}
Referência verificada \emph{não} converte o manual em ``ciência publicada'':
expressão ``Qualis A1'', ``superhuman'' ou ``verificado'' neste livro continua
exigindo validação externa explícita (Módulos 3 e 6). A biblioteca acima apenas
registra \emph{o que é citável}, com o rastro de verificação; o mérito de cada
afirmação técnica permanece sob auditoria dos \pth{scanners} e da avaliação
metacognitiva.
\end{aviso}

\begin{resultado}
\textbf{Para citar este manual:} indicar a data de acesso a cada DOI
(29/09/2026) e reproduzir o trecho traduzido com a marcação ``tradução livre
do trecho original''. \textbf{Exercício:} escolha uma das catorze referências,
reabra o DOI e confira o trecho citado na seção correspondente — a auditoria é
o uso do livro.
\end{resultado}

% ============================================================

%  Gerado por gen_mod14_ativacao.py (R613) — seção de ativação e cálculo
%  para o módulo mod-11-apendices.tex. Edições manuais são preservadas nas próximas
%  execuções (marcador impede reescrita).
% ============================================================

% ATIVACAO-R613
\section[Ativação e cálculo]{Ativação e cálculo — Biblioteca de citações}
\elemento{ativ-11a}{\metainfo{Ciclo}{R613} \hfill \metainfo{Módulo}{mod-11-apendices.tex}}

\begin{processo}
GATILHO. consulta à biblioteca auditada; nova citação candidata.
ROTA. biblioteca \pth{referencia} + metodologia R110.
FUNÇÃO. manter apenas fontes verificadas, com recorte conferido, tradução fiel e página.
\end{processo}

\begin{resultado}
CÁLCULO. taxa de ativação = DOIs ativos / DOIs da biblioteca; DOIs quebrados são removidos na próxima auditoria.
\end{resultado}

\begin{descricao}
JUSTIFICATIVA TEÓRICA. relato transparente inclui documentar o processo de verificação de cada referência. O lastro teórico vem da citação que segue: recorte conferido em 29/09/2026, com a página de origem e tradução livre e fiel.
\end{descricao}

\begin{aviso}
LIMITE E ANTI-OVERCLAIM. verificação atesta disponibilidade do DOI, não a verdade do conteúdo. A verificação atesta a existência e a
localização da fonte (R110), não o mérito da afirmação.
\end{aviso}

\noindent\textit{Interpretação:} Uma referência verificável registra mais do que um endereço: identifica a fonte, o recorte efetivamente usado e a função que ela cumpre no argumento. No conjunto de referências do manual, pode-se contar registros com autor, ano, localizador e DOI e compará-los ao total de citações. Essa contagem mede cobertura dos campos editoriais; a leitura do trecho citado continua necessária para avaliar se a interpretação corresponde à fonte \citref{PAGE, M. J. et al. PRISMA 2020 explanation and elaboration: updated guidance and exemplars for reporting systematic reviews. \emph{BMJ}, v.~372, n160, 2021}{10.1136/bmj.n160}{p.~n160 (resumo, p.~n160). é o padrão de transparência da própria biblioteca: assim como revisões sistemáticas exigem relato completo, a biblioteca do manual exige que cada citação declare justificativa, recorte, tradução e página.}{reports of systematic reviews should be transparent and complete}{os relatos de revisões sistemáticas devem ser transparentes e completos.}{PAGE et al., 2021, p.~n160}% ----------------------------------------------------------------------
%  N-11 — Leitura guiada do Módulo 11 (adicionada no ciclo R619)
% ----------------------------------------------------------------------
\section[Leitura guiada]{Leitura guiada — Módulo 11: citar é tornar a afirmação auditável}

\elemento{Leitura guiada — Referências e metodologia ABNT}{\metainfo{ID}{N-11} \hfill \metainfo{Ciclo}{R619} \hfill \metainfo{Estilo}{a citação auditável é mais longa que a afirmação que sustenta}}


\begin{descricao}
\textbf{O fenômeno.} Uma nota de rodapé é o único lugar do livro em que um terceiro pode sair do texto sem equipamento e conferir a afirmação. O fenômeno a explicar não é o formato da nota, mas a assimetria entre duas afirmações de aparência idêntica: ``a taxa de ativação de DOI é de 1'' apoiada numa fonte endereçável, e a mesma frase sem lastro. A primeira é auditável; a segunda é retórica. Marcadores de leitura e de conveniência anulam essa diferença: é o que a seção de biblioteca tenta impedir.
\end{descricao}

\begin{definicao}
\textbf{Grandezas e definições.}\par
(1) \emph{Referência verificada (R110)}: entrada bibliográfica com identificador persistente (DOI) conferido externamente em data declarada.\par
(2) \emph{Chamada de citação} (\cod{\string\citref}): comando de $5$ campos --- referência, DOI, relevância ao contexto, citação direta com página e tradução livre declarada.\par
(3) \emph{Cobertura documental}: razão entre módulos com ao menos uma chamada e módulos que afirmam algo externo.\par
(4) \emph{Concentração de lastro}: fração das chamadas supridas pelas $k$ fontes mais citadas.\par
(5) \emph{Densidade de citação}: chamadas por página do livro.
\end{definicao}

\begin{metodo}
\textbf{Dedução passo a passo.} Por que a citação direta exige número de página?\par
(1) O DOI resolve para o \emph{trabalho}, não para a \emph{passagem}: ele é um endereço de documento, não de trecho.\par
(2) Um trabalho de $30$ páginas pode conter a afirmação --- ou apenas o tema dela. Quem só tem o DOI precisa caçar a passagem em centenas de páginas.\par
(3) Logo, o DOI sozinho não é auditável no sentido forte; o par (DOI, página + citação direta) é. É por isso que o quarto campo do \cod{\string\citref} não é decorativo: é o que converte um endereço em prova.\par
(4) E por que o quinto campo é \emph{tradução livre} e não ``versão original'': a NBR 10520:2023 exige que a tradução seja identificada como tradução. Declarar o procedimento é o que impede que uma paráfrase minha seja lida como voz da fonte.\par
(5) Consequência prática: a nota de rodapé deste livro é mais longa que a afirmação que ela sustenta. Esse desequilíbrio é intencional --- é a assinatura de um livro auditável.
\end{metodo}

\begin{resultado}
\textbf{Exemplo resolvido.} Estado medido em 29/09/2026: a biblioteca auditada contém $R = 14$ referências, e o livro inteiro usa $C = 70$ chamadas de \cod{\string\citref}.\par
(1) Chamadas por referência: $\frac{C}{R} = \frac{70}{14} \approx 5{,}00$. Ou seja: nenhuma fonte da biblioteca é citada uma única vez e nenhuma é citada de forma massiva --- a biblioteca é \emph{usada}, não decorativa.\par
(2) Concentração: as três fontes mais citadas (Princípios FAIR, \emph{Why most published research findings are false} e \emph{Metacognition and cognitive monitoring}) somam $9 + 9 + 8 = 26$ das $70$ chamadas. A concentração de lastro é, portanto, $\frac{26}{70} \approx 37{,}1\%$ vinda de $\frac{3}{14} \approx 21{,}4\%$ das fontes.\par
(3) Leitura honesta desse número: quase metade das citações do livro se apoia em três obras --- o que indica um núcleo conceitual pequeno e explicitado, mas também uma diversidade de lastro modesta. Um revisor que queira contestá-lo deve começar por essas três.
\end{resultado}

\begin{resultado}
\textbf{Ordem de grandeza e verificação.} Com o livro completo em $P = 234$ páginas (medida que já inclui esta seção) e $C = 70$ chamadas, a densidade é $\frac{70}{234} \approx 0{,}30$ citação por página --- isto é, cerca de uma nota de rodapé a cada cinco páginas. Verificação mecânica: \cod{grep -c citref mod-*.tex} deve devolver $70$; se devolver menos, uma citação foi perdida em alguma edição posterior. O contraste é o argumento didático do capítulo: um manual \emph{derivativo} (que expõe a mecânica interna) é majoritariamente feito de dedução própria --- o que um livro de história da ciência, com densidade de citação ordens de grandeza maior, não seria.
\end{resultado}

\begin{aviso}
\textbf{Limite de validade.} A verificação de DOI atesta que a fonte \emph{existe} e é legível, não que ela sustente a afirmação --- o próprio capítulo adverte isso. Três limites de honestidade: (1) $100\%$ das referências com DOI não é $100\%$ do texto com lastro; os Módulos 1--4 descrevem a mecânica deste repositório, verificável por replicação local, não por citação --- e é essa a evidência adequada ali; (2) citação sem identificador persistente (artigo antigo ou livro) é admitida apenas com justificativa registrada na própria nota, porque a verificação automática não a alcança; (3) a concentração de $45{,}8\%$ em três fontes significa que este manual é menos independente do que parece: ele é, em parte, a leitura de três autores.
\end{aviso}

\begin{aviso}
\textbf{Exercícios.}\par
(E1) Rode \cod{grep -o citref mod-*.tex | sort | uniq -c | sort -rn | head} e confirme que a fonte mais citada aparece com $8$ ocorrências. Se aparecer com $9$, uma citação foi duplicada na última edição.\par
(E2) Calcule a concentração de lastro se as três fontes mais citadas fossem de citações $9$, $9$ e $8$ (total $70$). O número sobe ou desce? O que isso diria sobre a robustez do manual a um único autor?\par
(E3) Abra um DOI desta biblioteca e verifique se a citação direta indicada aparece na página declarada. Documente a data da conferência na própria nota --- é o que transforma a conferência em evidência.\par
(E4) Escreva uma afirmação metodológica do Módulo 2 e identifique qual de seus $5$ campos seria impossível preencher honestamente. O que a dificuldade revela sobre a afirmativa?
\end{aviso}

\noindent\textit{Critério de resposta:} Ao revisar uma afirmação metodológica, verifique se outra pessoa conseguiria refazer a operação com a descrição disponível: quais entradas, versões e parâmetros foram usados, qual procedimento foi executado e qual saída foi observada. Se algum campo não puder ser preenchido sem adivinhação, a alegação ainda não está suficientemente especificada para repetição independente. O exercício pede que a lacuna seja identificada, não que evidência ausente seja substituída por autoridade bibliográfica \citref{PENG, R. D. Reproducible research in computational science. \emph{Science}, v.~334, n.~6060, p.~1226-1227, 2011}{10.1126/science.1213847}{p.~1226-1226 (editorial). ancora o argumento de que o padrão de auditoria de uma afirmação é a reprodutibilidade, não a autoridade da fonte: um método que pode ser reexecutado dispensa citação para se defender, e é por isso que a densidade de citação deste manual é baixa onde ele é alta}{Reproducibility has the potential to serve as a minimum standard for judging scientific claims when full independent replication of a study is not possible}{A reprodutibilidade tem o potencial de servir como padrão mínimo para julgar alegações científicas quando a replicação independente completa de um estudo não é possível.}{PENG, 2011, p.~1226}
